\begin{titlepage}
\centering
\vspace*{2cm}
{\Huge\bfseries AUTHBC\par}
\vspace{0.6cm}
{\LARGE Feasibility Boundaries for Authenticated UAV Telemetry\par}
\vspace{0.3cm}
{\large An Exclusion Bound, a Capacity Envelope, and Hardware Validation\par}
\vspace{2.5cm}
{\Large Mohamed A. Farouk\par}
\vspace{1.5cm}
{\large A thesis submitted for the degree of\par}
\vspace{0.3cm}
{\large [DEGREE --- TO BE COMPLETED]\par}
\vspace{1.5cm}
{\large [DEPARTMENT, INSTITUTION --- TO BE COMPLETED]\par}
\vspace{2cm}
{\large \today\par}
\vfill
{\small\itshape Draft. Chapter completeness is declared in \texttt{thesis/STATUS.md}.\par}
\end{titlepage}

\chapter*{Abstract}
\addcontentsline{toc}{chapter}{Abstract}
\markboth{Abstract}{Abstract}

Flying ad-hoc networks (FANETs) stream telemetry records so small that a cryptographic signature
($64$--$96$\,B) rivals or exceeds the payload itself ($45$--$190$\,B). Once the payload has been
compressed, authentication --- not data --- becomes the dominant on-air cost, and the operational
question stops being \emph{how many bytes does a scheme save} and becomes \emph{which
configurations can run at all}.

This thesis answers the second question. It searches four design axes jointly --- record
\emph{encoding}, authentication \emph{placement}, signature \emph{scheme} and \emph{batch} size ---
against a byte model, an NS-3-validated IEEE 802.11 broadcast channel model, and energy measured
on Raspberry~Pi~4 hardware. Three results follow, ordered by how durable they are.

\textbf{An exclusion.} A link whose payload cannot hold one frame header, one signature and one
record admits no per-frame-verifiable telemetry at any encoding or batch size, and packet loss
forecloses fragmentation as an escape, because a unit spanning $n$ frames verifies with probability
$(1-p)^n$. On the EU868 LoRa band this removes \emph{three} of seven data rates unconditionally: at
DR0--DR2 a $64$\,B signature alone overflows a $51$\,B payload, so the exclusion survives a
zero-byte header and a one-byte record. A fourth rate, DR3, misses by six bytes against a measured
frame header of $44$\,B --- and is recovered by an integer-keyed header profile, so the boundary is
reported \emph{constructively}: it names not only where authentication cannot fit, but what would
have to change.

\textbf{A capacity envelope.} Joining the byte model to the channel model yields the largest
neighbourhood each configuration can sustain. The co-designed configuration carries
$1.9\times$--$3.2\times$ the neighbourhood of an inline-signing baseline, reported across both the
saturation and the verifiability threshold rather than at whichever flatters the result.

\textbf{Bytes, and what they cost.} The adopted configuration puts $58.7\%$ fewer bytes on air
($72.0$\,B per record) at an operating point compliant with 3GPP TS~22.125, against a success
criterion committed to version control two days before the result existed.

A factorial ablation narrows rather than inflates the co-design claim: only placement and batching
interact, by exactly $\ga(1-1/b)$; encoding contributes additively and the signature-scheme axis is
byte-neutral in this regime. AUTHBC introduces no new cryptographic primitive. Its contribution is
a reproducible co-design, its quantified trade-offs, an honest account of the boundary beyond which
no choice of existing cryptography helps --- and a methodological record of the defects found while
establishing it.

\chapter*{Declaration}
\addcontentsline{toc}{chapter}{Declaration}
\markboth{Declaration}{Declaration}

\needswork{TO BE COMPLETED by the candidate. Institutional wording required.}

\vspace{1em}
\noindent\textbf{Use of generative AI.} \needswork{The candidate must state this explicitly and
accurately; the repository history records the extent of agent involvement in implementation,
experiment execution and drafting. Institutional policy governs the required wording.}

\chapter*{Acknowledgements}
\addcontentsline{toc}{chapter}{Acknowledgements}
\markboth{Acknowledgements}{Acknowledgements}

\needswork{TO BE COMPLETED.}

\chapter*{Nomenclature}
\addcontentsline{toc}{chapter}{Nomenclature}
\markboth{Nomenclature}{Nomenclature}

Symbols follow \texttt{docs/01\_SYSTEM\_MODEL\_ARCHITECTURE.md}~\S2, which is the single source of
truth for notation across the repository, the papers and this thesis.

\begin{center}
\begin{tabular}{ll l}
\toprule
Symbol & Meaning & Default \\
\midrule
$s$        & encoded record payload bytes & measured (E1) \\
$\ga$      & authentication object bytes (scheme $\sigma$) & 64 / 64 / 96 \\
$\Hf$      & ledger frame header bytes & \textbf{38--44, measured} \\
$M$        & application MTU budget & 1500 (802.11) \\
$b$        & records per frame (batch) & decision variable \\
$p$        & frame loss probability & $\{0.02, 0.05, 0.10\}$ \\
$\varepsilon$ & verifiability slack; require $V \geq 1-\varepsilon$ & $=p$ \\
$\Lambda_i$ & per-UAV record rate (rec/s) & 50 (adopted) \\
$\Lambda$   & aggregate rate a receiver verifies & $\Lambda_i \cdot N_{\text{local}}$ \\
$N_{\text{local}}$ & neighbours in the collision domain & reported as a curve \\
$U$        & channel utilisation, offered / deliverable & admission metric \\
$V$        & $P(\text{record verifiable at receiver})$ & constraint $\geq 1-\varepsilon$ \\
$R$        & PHY data rate & \SI{6}{\mega\bit\per\second} \\
$D(b)$     & freshness of the oldest record in a batch & $\leq \Dmax$ \\
$\smax$    & largest record a single frame admits & $M - \Hf - \ga$ \\
$\Nmax$    & largest sustainable neighbourhood & derived \\
\bottomrule
\end{tabular}
\end{center}

\vspace{1em}
\noindent\textbf{A note on $\Hf$.} It is reported as a \emph{range}, not a constant. Canonical CBOR
encodes small integers in one byte and large ones in five, so the measured frame header moves
between $38$ and $44$\,B with the magnitudes of the sender identifier and sequence number. This
matters because one of the thesis's own results depends on it; \Cref{ch:lowrate} states where.
